In recent years, we have started to see a surge of LLM or large language models being used in everyday life as a powerful personal helper. This approach helps millions with every day tasks, but it also starts replacing many automated task to be more AI driven. LLMs are at the core of many recommendation pipelines for retrieval, re-ranking and explanation generation. It thus shows a clear shift in the way we evaluate these recommender systems, as they are more prone to be manipulated by harmful actors. This proposal presents this core issue through our subject, \emph{Red-Teaming Large Language Model-Powered Recommendation Systems}. With the use of my experimental project RAG Poison Lab, we will be able to track poisoning signals in every step of the retrieval process until the final recommendation output\cite{ragpoisonlab_repo} using real life metrics to get a measurable outputs, HR@K, NDCG@K, MRR@K, ASR@K, candidate overlap, and target-rank lift. We are hoping to answer the following questions: where is the earliest step in which an actor can attack the system, and how much can they manipulate the final output. Our plan will run in a controlled environment where we will compare the baseline recommender results to the attacked ones simulating the real experience of an attacked RAG system.